% FAC-EN-U0002
% French source: fac_body.tex lines 116--162.
% Authority: PDF physical pp. 5--6; printed pp. 200--201 (no. 3 complete).

\medskip
\phantomsection\label{fac:I:1:no:3}
\textbf{3. Construction of sheaves.} Suppose that, for every open set $U\subset X$, an abelian group $\mathfrak{F}_U$ is given and, for every pair of open sets $U\subset V$, a homomorphism $\varphi_U^V:\mathfrak{F}_V\to\mathfrak{F}_U$ is given, in such a way that the transitivity condition $\varphi_U^V\circ\varphi_V^W=\varphi_U^W$ holds whenever $U\subset V\subset W$.

The collection $(\mathfrak{F}_U,\varphi_U^V)$ then defines a sheaf $\mathfrak{F}$ as follows:

\begin{enumerate}[label=(\alph*),leftmargin=2.5em]
\item\label{fac:I:1:no:3:item:a} Put $\mathfrak{F}_x=\varinjlim\mathfrak{F}_U$, the direct limit over the directed set of open neighborhoods $U$ of $x$. Thus, if $x$ belongs to the open set $U$, there is a canonical homomorphism $\varphi_x^U:\mathfrak{F}_U\to\mathfrak{F}_x$.

\item\label{fac:I:1:no:3:item:b} If $t\in\mathfrak{F}_U$, let $[t,U]$ denote the set of the elements $\varphi_x^U(t)$ as $x$ ranges over $U$. Thus $[t,U]\subset\mathfrak{F}$, and we give $\mathfrak{F}$ the topology generated by the sets $[t,U]$. Consequently, an element $f\in\mathfrak{F}_x$ has as a neighborhood basis in $\mathfrak{F}$ the sets $[t,U]$ for which $x\in U$ and $\varphi_x^U(t)=f$.
\end{enumerate}

It follows immediately that the data \hyperref[fac:I:1:no:3:item:a]{(a)} and \hyperref[fac:I:1:no:3:item:b]{(b)} satisfy axioms \hyperref[fac:I:1:no:1:axiom:I]{(I)} and \hyperref[fac:I:1:no:1:axiom:II]{(II)}, that is, that $\mathfrak{F}$ is indeed a sheaf. We shall call it the sheaf \textit{defined by the system} $(\mathfrak{F}_U,\varphi_U^V)$.

If $t\in\mathfrak{F}_U$, the map $x\mapsto\varphi_x^U(t)$ is a section of $\mathfrak{F}$ over $U$; hence there is a canonical homomorphism $\iota:\mathfrak{F}_U\to\Gamma(U,\mathfrak{F})$.

\medskip
\phantomsection\label{fac:I:1:no:3:prop:1}
\textsc{Proposition 1.} \textit{The homomorphism $\iota:\mathfrak{F}_U\to\Gamma(U,\mathfrak{F})$ is injective\footnote{Recall (cf. [1]) that a map $f:E\to E'$ is called \textit{injective} if $f(e_1)=f(e_2)$ implies $e_1=e_2$, \textit{surjective} if $f(E)=E'$, and \textit{bijective} if it is both injective and surjective. An injective (respectively surjective, bijective) map is called an injection (respectively a surjection, a bijection).\label{fac:footnote:inj-surj}} if and only if the following condition holds:}

\textit{If $t\in\mathfrak{F}_U$ and there exists an open covering $\{U_i\}$ of $U$ such that $\varphi_{U_i}^U(t)=0$ for every $i$, then $t=0$.}

If $t\in\mathfrak{F}_U$ satisfies the preceding condition, then
\begin{equation*}
\varphi_x^U(t)=\varphi_x^{U_i}\circ\varphi_{U_i}^U(t)=0
\qquad\text{if }x\in U_i,
\end{equation*}
which means that $\iota(t)=0$. Conversely, suppose that $\iota(t)=0$ for some $t\in\mathfrak{F}_U$. Since $\varphi_x^U(t)=0$ for $x\in U$, the definition of the direct limit gives an open neighborhood $U(x)$ of $x$ such that $\varphi_{U(x)}^U(t)=0$. The sets $U(x)$ then form an open covering of $U$ satisfying the condition in the statement.

\medskip
\phantomsection\label{fac:I:1:no:3:prop:2}
\textsc{Proposition 2.} \textit{Let $U$ be an open subset of $X$, and suppose that $\iota:\mathfrak{F}_V\to\Gamma(V,\mathfrak{F})$ is injective for every open set $V\subset U$. The homomorphism $\iota:\mathfrak{F}_U\to\Gamma(U,\mathfrak{F})$ is surjective\textsuperscript{\ref{fac:footnote:inj-surj}} (hence bijective) if and only if the following condition holds:}

\textit{For every open covering $\{U_i\}$ of $U$ and every family $\{t_i\}$, where $t_i\in\mathfrak{F}_{U_i}$ and}
\begin{equation*}
\varphi_{U_i\cap U_j}^{U_i}(t_i)=\varphi_{U_i\cap U_j}^{U_j}(t_j)
\end{equation*}
\textit{for every pair $(i,j)$, there exists $t\in\mathfrak{F}_U$ such that $\varphi_{U_i}^{U}(t)=t_i$ for every $i$.}

The condition is necessary. Each $t_i$ defines a section $s_i=\iota(t_i)$ over $U_i$, and $s_i=s_j$ over $U_i\cap U_j$. Hence there exists a section $s$ over $U$ which agrees with $s_i$ on $U_i$ for every $i$. If $\iota:\mathfrak{F}_U\to\Gamma(U,\mathfrak{F})$ is surjective, there exists $t\in\mathfrak{F}_U$ such that $\iota(t)=s$. Put $t_i'=\varphi_{U_i}^{U}(t)$. The section defined by $t_i'$ over $U_i$ is precisely $s_i$; hence $\iota(t_i-t_i')=0$, which implies $t_i=t_i'$ because $\iota$ is assumed to be injective.

The condition is sufficient. If $s$ is a section of $\mathfrak{F}$ over $U$, there exist an open covering $\{U_i\}$ of $U$ and elements $t_i\in\mathfrak{F}_{U_i}$ such that $\iota(t_i)$ is the restriction of $s$ to $U_i$. It follows that $\varphi_{U_i\cap U_j}^{U_i}(t_i)$ and $\varphi_{U_i\cap U_j}^{U_j}(t_j)$ define the same section over $U_i\cap U_j$ and are therefore equal, by the hypothesis on $\iota$. If $t\in\mathfrak{F}_U$ satisfies $\varphi_{U_i}^{U}(t)=t_i$, then $\iota(t)$ agrees with $s$ on every $U_i$, and hence on $U$. \hfill$\square$

\medskip
\phantomsection\label{fac:I:1:no:3:prop:3}
\textsc{Proposition 3.} \textit{If $\mathfrak{F}$ is a sheaf of abelian groups on $X$, the sheaf defined by the system $(\Gamma(U,\mathfrak{F}),\rho_U^V)$ is canonically isomorphic to $\mathfrak{F}$.}

This follows immediately from the properties of sections stated in no.~\hyperref[fac:I:1:no:2]{2}.

\hyperref[fac:I:1:no:3:prop:3]{Proposition 3} shows that every sheaf can be defined by a suitable system $(\mathfrak{F}_U,\varphi_U^V)$. Note that different systems may define the \textit{same} sheaf $\mathfrak{F}$. Nevertheless, if the systems $(\mathfrak{F}_U,\varphi_U^V)$ are required to satisfy the conditions of \hyperref[fac:I:1:no:3:prop:1]{Propositions 1} and \hyperref[fac:I:1:no:3:prop:2]{2}, there is, up to isomorphism, only one possible system: the one formed by $(\Gamma(U,\mathfrak{F}),\rho_U^V)$.

\textsc{Example.} Let $G$ be an abelian group, and take $\mathfrak{F}_U$ to be the set of functions on $U$ with values in $G$. Define $\varphi_U^V:\mathfrak{F}_V\to\mathfrak{F}_U$ by restricting a function. This gives a system $(\mathfrak{F}_U,\rho_U^V)$ and hence a sheaf $\mathfrak{F}$, called the \textit{sheaf of germs of functions} with values in $G$. It is immediate that the system $(\mathfrak{F}_U,\varphi_U^V)$ satisfies the conditions of \hyperref[fac:I:1:no:3:prop:1]{Propositions 1} and \hyperref[fac:I:1:no:3:prop:2]{2}; consequently, the sections of $\mathfrak{F}$ over an open set $U$ may be identified with the elements of $\mathfrak{F}_U$.
